Despite progress on short-horizon tasks, large language model (LLM) agents struggle to sustain reliable decision-making in realistic environments, particularly in real-world economic settings. Yet where their decision-making breaks down remains poorly understood. We introduce RetailBench, a controlled, data-grounded benchmark for long-horizon autonomous business operations in realistic settings. We evaluate seven LLMs with multiple agent scaffolds and conduct fine-grained analyses of their decision-making processes. Results over a 180-day evaluation horizon show that most agents substantially underperform traditional operations research methods. 
 Even the best-performing agent achieves only 30.1\% of the final net worth attained by a traditional operations research baseline and 19.6\% of that achieved by a privileged reference policy.Our analysis reveals that agents gather evidence effectively but struggle to make high-quality decisions under complex dependencies and respond promptly to operational anomalies. Despite achieving 99.4\% fully evidenced actions, the best-performing agent attains only a 54.4\% decision-making score and a 27.6\% three-day anomaly response rate.
These findings highlight the persistent challenges of maintaining reliable agent performance over long horizons in realistic environments, positioning RetailBench as a testbed for evaluating and diagnosing long-horizon agent capabilities.